\documentclass[11pt]{article}
\usepackage{amsmath, amssymb, amsfonts, amsthm, geometry, bm, booktabs}
\usepackage[hidelinks]{hyperref}
\geometry{margin=1in}

\title{Discrete Information Geometry of Prime Triplets: \\ The Modulo-30 Arithmetic Lattice and Variance Suppression}
\author{Rob Miller}
\date{June 2026}

\newtheorem{theorem}{Theorem}
\newtheorem{lemma}{Lemma}
\newtheorem{corollary}{Corollary}
\newtheorem{conjecture}{Conjecture}
\newtheorem{remark}{Remark}
\newtheorem{observation}{Observation}

\begin{document}
\maketitle

\begin{abstract}
We extend the discrete information-geometric framework of prime constellations to prime triplets $(p, p+2, p+6)$ and $(p, p+4, p+6)$. While twin primes are structurally constrained to a modulo-6 arithmetic lattice, prime triplets are governed by a strictly tighter modulo-30 lattice ($\mathcal{L}_{30}$). We formalize the discrete information geometry for $\mathcal{L}_{30}$, deriving the Fisher information metric $I(\lambda) = 225/[\lambda^4 \sinh^2(15/\lambda)]$, and demonstrate that imposing continuous exponential geometries on triplet gaps mischaracterizes the discrete structure at computable scales. 

Complementing the theoretical lattice structure, we present an empirical equidistribution census for prime triplets across triplet-admissible residue classes modulo $P_n$. Building on the sub-Poisson variance suppression observed in twin primes, the triplet census reveals a pronounced structural anti-clumping effect at small primorial scales, with the Coefficient of Variation strictly below the Poisson baseline. This work constitutes Paper 3 in the post-v4.3 Recursive Admissibility Prime Framework (RAPF) trilogy, demonstrating that sub-Poisson variance suppression is a finite-scale boundary effect that systematically attenuates toward Poisson randomness as $n \to \infty$, for both $k=2$ and $k=3$ constellations.
\end{abstract}

\noindent \textbf{MSC 2020:} 11N05, 11N36, 94A15

\section{Introduction}

The Hardy--Littlewood $k$-tuples conjecture provides the foundational heuristic for the asymptotic density of prime constellations. For prime triplets of the form $\mathcal{H}_1 = (0, 2, 6)$ and $\mathcal{H}_2 = (0, 4, 6)$, the conjectured density scales as $\lambda(X) \sim (\log X)^3 / \mathfrak{S}_3$, where $\mathfrak{S}_3$ is the triplet singular series constant. 

As established in Paper 1 of this program, comparing finite-scale empirical gap data to continuous asymptotic predictions creates artificial discrepancies if the underlying arithmetic constraints are ignored. Twin prime gaps do not exist in $\mathbb{R}^+$; they are strictly confined to a modulo-6 lattice ($\mathcal{L}_6 = \{6k : k \in \mathbb{N}\}$). Imposing a smooth density on a discrete geometry fails to capture the structural information of the lattice.

Prime triplets face an even more severe geometric constraint. For $p, p+2,$ and $p+6$ to all be prime, $p$ must satisfy simultaneous congruences modulo 2, 3, and 5. This restricts the anchor prime $p$ to a specific subset of residue classes modulo 30 (e.g., $p \equiv 11 \pmod{30}$ for $\mathcal{H}_1$, $p \equiv 7 \pmod{30}$ for $\mathcal{H}_2$). Consequently, the gap between consecutive prime triplets is strictly a multiple of 30.

This paper extends the discrete information-geometric framework to the modulo-30 lattice ($\mathcal{L}_{30}$) and presents the corresponding equidistribution census. We demonstrate that the sub-Poisson variance suppression observed for twin primes (Paper 2) is a generalized feature of admissible classes at small primorial scales, and that this suppression attenuates toward Poisson as $n \to \infty$ for both $k=2$ and $k=3$ constellations.

\section{Discrete Information Geometry on $\mathcal{L}_{30}$}

\subsection{The Modulo-30 Lattice PMF}

Triplet-to-triplet gaps $g$ take values strictly in $\mathcal{L}_{30} = \{30k : k = 1, 2, 3, \ldots\}$. We model the lattice index $k = g/30$ using a discrete geometric probability mass function with an exponential envelope, scaled by the dynamic Hardy-Littlewood parameter $\lambda(X)$:
\begin{equation}
    P(k; \lambda) = \left(e^{30/\lambda} - 1\right) e^{-30k/\lambda}, \qquad k \in \{1, 2, 3, \ldots\}.
\end{equation}
This distribution is a geometric PMF with success parameter $p = 1 - e^{-30/\lambda}$. The mean and variance of the index $k$ are:
\begin{equation}
    \mathbb{E}[k] = \frac{e^{30/\lambda}}{e^{30/\lambda} - 1}, \qquad \text{Var}(k) = \frac{e^{-30/\lambda}}{(1 - e^{-30/\lambda})^2}.
\end{equation}

\subsection{Exact Fisher Information for Triplets}

The log-likelihood is $\log P(k; \lambda) = \log(e^{30/\lambda} - 1) - 30k/\lambda$, with score function:
\begin{equation}
    \frac{\partial}{\partial \lambda} \log P(k; \lambda) = \frac{30}{\lambda^2}\left(k - \mathbb{E}[k]\right).
\end{equation}

The Fisher information metric is the variance of the score:
\begin{align}
    I(\lambda) &= \mathbb{E}\left[\left(\frac{\partial}{\partial \lambda} \log P(k; \lambda)\right)^2\right] = \frac{900}{\lambda^4}\,\text{Var}(k) \nonumber\\
    &= \frac{900}{\lambda^4} \cdot \frac{e^{-30/\lambda}}{(1 - e^{-30/\lambda})^2}.
\end{align}

Using the hyperbolic identity $4\sinh^2(x) = e^{2x} - 2 + e^{-2x}$ with $x = 15/\lambda$, we obtain the exact closed-form metric:

\begin{theorem}[Discrete Triplet Fisher Information]
For the geometric PMF on the modulo-30 lattice $\mathcal{L}_{30}$ with dynamic scale parameter $\lambda > 0$, the Fisher information is:
\begin{equation}
    I(\lambda) = \frac{225}{\lambda^4 \sinh^2(15/\lambda)}.
\end{equation}
\end{theorem}

\begin{remark}[Continuous Limit]
As $\lambda \to \infty$, the small-angle approximation $\sinh(15/\lambda) \approx 15/\lambda$ yields $I(\lambda) \to 225/[\lambda^4 (225/\lambda^2)] = 1/\lambda^2$. Thus, the discrete modulo-30 geometry flawlessly recovers the continuous exponential information limit. However, at scales $\lambda \sim O(30)$, the $\sinh$ formulation is geometrically required because the lattice spacing is not negligible relative to $\lambda$.
\end{remark}

\section{Admissible Classes and Lattice Structure}

\subsection{Triplet-Admissible Residues}

A residue class $a \pmod{P_n}$ is triplet-admissible for $\mathcal{H}_1 = (0, 2, 6)$ if $a, a+2,$ and $a+6$ are all coprime to $P_n$. The count factors per prime:
\begin{itemize}
  \item Modulo 2: $a$ must be odd (1 choice).
  \item Modulo 3: $a \equiv 0$ is forbidden, and $(a+2) \equiv 0 \implies a \equiv 1$ is also forbidden, leaving $a \equiv 2$ (1 choice).
  \item Modulo 5: residues $a \equiv 0, 3, 4 \pmod 5$ are forbidden ($(a+2)$ and $(a+6)$ cover 2 of the remaining), leaving 2 choices.
  \item For $p \geq 7$: the three residues $a, a+2, a+6 \pmod p$ are distinct, so exactly 3 of the $p$ classes are forbidden, leaving $p-3$ choices.
\end{itemize}
The total count of triplet-admissible classes is therefore:
\begin{equation}
    A_{n,3} = 2 \cdot \prod_{7 \le p \le p_n} (p-3).
\end{equation}

\begin{table}[h]
\centering
\caption{Triplet-admissible class counts for small $n$}
\label{tab:triplet_counts}
\begin{tabular}{ccccc}
\toprule
$n$ & $P_n$ & Triplet classes $A_{n,3}$ & Twin classes $A_n$ & Ratio $A_{n,3}/A_n$ \\
\midrule
3 & 30 & 2 & 3 & 0.67 \\
4 & 210 & 8 & 15 & 0.53 \\
5 & 2{,}310 & 64 & 135 & 0.47 \\
6 & 30{,}030 & 640 & 1{,}485 & 0.43 \\
7 & 510{,}510 & 8{,}960 & 22{,}275 & 0.40 \\
\bottomrule
\end{tabular}

\vspace{0.5em}
\noindent{\footnotesize Class counts $A_{n,3}$ shown per pattern. Combined census (Table~\ref{tab:triplet_census}) uses $2 \times A_{n,3}$ classes to account for both $\mathcal{H}_1 = (0, 2, 6)$ and $\mathcal{H}_2 = (0, 4, 6)$.
}
\end{table}

The ratio $A_{n,3}/A_n$ decreases with $n$ because for $p \geq 7$, the per-prime growth factor for triplets is $(p-3)/2$ while for twins it is $(p-2)$. The ratio $(p-3)/(2(p-2)) < 1$ for all $p \geq 7$, meaning the combinatorial space for $k=3$ constellations thins significantly faster than for twins.

\section{Equidistribution Census of Prime Triplets}

We executed a parallelized segmented sieve to capture triplet occurrences across all triplet-admissible classes up to primorial $P_n$. The census tracks both triplet patterns $\mathcal{H}_1 = (0, 2, 6)$ and $\mathcal{H}_2 = (0, 4, 6)$ independently, then reports combined statistics. For each scale $n$, the sieve covers all classes across the full range $[P_n, P_n^2]$, ensuring complete coverage with no sampling.

\begin{table}[h]
\centering
\caption{Equidistribution of Prime Triplets across admissible classes mod $P_n$}
\label{tab:triplet_census}
\begin{tabular}{ccccccccc}
\toprule
$n$ & $P_n$ & $2A_{n,3}$ & Total Triplets & Mean ($\mu$) & CV & CV$_{\rm Poisson}$ & Ratio & Empty \\
\midrule
3 & 30 & 4 & 25 & 6.3 & 23.66\% & 40.00\% & 0.592 & 0 \\
4 & 210 & 16 & 284 & 17.8 & 12.83\% & 23.74\% & 0.541 & 0 \\
5 & 2{,}310 & 128 & 10{,}381 & 81.1 & 9.50\% & 11.10\% & 0.856 & 0 \\
6 & 30{,}030 & 1{,}280 & 695{,}318 & 543.2 & 3.98\% & 4.29\% & 0.928 & 0 \\
\textbf{7} & \textbf{510{,}510} & \textbf{17{,}920} & \textbf{93{,}169{,}155} & \textbf{5{,}199.2} & \textbf{1.30\%} & \textbf{1.39\%} & 0.935 & \textbf{0} \\
\bottomrule
\end{tabular}

\vspace{0.5em}
\noindent{\footnotesize CV$_{\rm Poisson} = 100/\sqrt{\mu}$ represents the independent Poisson baseline. Ratio $=$ CV/CV$_{\rm Poisson}$ quantifies sub-Poisson suppression (1.0 = Poisson, $<1.0$ = sub-Poisson). Combined classes = $2 \times A_{n,3}$ (both patterns). Total triplets combines both $\mathcal{H}_1 = (0, 2, 6)$ and $\mathcal{H}_2 = (0, 4, 6)$. Ranges: $n=3$ covers $[30, 900]$, $n=4$ covers $[210, 44{,}100]$, $n=5$ covers $[2{,}310, 5{,}336{,}100]$, $n=6$ covers $[30{,}030, 9.02 \times 10^8]$, $n=7$ covers $[510{,}510, 2.61 \times 10^{11}]$.
}
\end{table}

The census reveals three key findings. First, the Coefficient of Variation decreases monotonically from 23.66\% at $n=3$ to 1.30\% at $n=7$, demonstrating that triplet equidistribution tightens as the lattice becomes more restrictive. Second, at every scale studied, the observed CV falls strictly below the independent Poisson baseline CV$_{\rm Poisson}$, confirming sub-Poisson variance suppression for triplets---a phenomenon first identified for twin primes at $n=8$ (CV = 0.15\%, Paper 2). Third, with zero empty classes across all 17{,}920 triplet-admissible classes at $n=7$, the joint class occupancy conjecture extends to $k=3$ constellations.

However, the raw CV drop does not imply intensifying suppression: the CV/CV$_{\rm Poisson}$ ratio (added as the new ``Ratio'' column) reveals the true trajectory of deviation from Poisson. At small scales ($n=3,4$) the ratio is furthest from $1.0$, indicating the strongest sub-Poisson effect (41--46\% suppression). As $n$ grows, the ratio steadily approaches unity---0.935 at $n=7$ (6.5\% suppression). This is the correct reading: sub-Poisson suppression is a finite-scale boundary effect that systematically relaxes toward Poisson as the number of admissible classes grows large. The phenomenon is consistent with a multinomial constraint operating at small $n$, where a fixed total distributed over a handful of bins produces measurable variance reduction relative to an idealized infinite-population Poisson. As $n \to \infty$ and the number of classes increases into the thousands, this effect attenuates and the ratio converges toward $1.0$.

\begin{table}[h]
\centering
\caption{Individual pattern census at $n=7$ for Pattern A $(p, p+2, p+6)$ and Pattern B $(p, p+4, p+6)$}
\label{tab:triplet_split}
\begin{tabular}{lcccccc}
\toprule
Pattern & Total & Classes & Mean & CV & Min & Max \\
\midrule
A $(0, 2, 6)$ & 46{,}591{,}342 & 8{,}960 & 5{,}199.9 & 1.2952\% & 4{,}950 & 5{,}433 \\
B $(0, 4, 6)$ & 46{,}577{,}813 & 8{,}960 & 5{,}198.4 & 1.3141\% & 4{,}966 & 5{,}473 \\
\bottomrule
\end{tabular}

\vspace{0.5em}
\noindent{\footnotesize The mirror-pattern CV split is just $0.02\%$, demonstrating that the geometric pressure is uniform across chiral configurations. The lattice does not distinguish left-leaning from right-leaning triplet topologies.
}
\end{table}

The mirror symmetry between Pattern A and Pattern B is a notable structural feature. Over 46 million triplet pairs per pattern, the two chiral configurations exhibit nearly identical dispersion (CV difference of $0.02\%$), identical class occupancy (zero empties), and total counts that differ by only $0.03\%$. This statistical equilibrium demonstrates that the variance suppression mechanism acts uniformly on both orientations of the triplet lattice, consistent with large-sample stability across chiral configurations.

Combined with the twin prime census trajectory (Paper 2: CV from 2.44\% at $n=3$ to 0.15\% at $n=8$), we now have a unified picture: sub-Poisson variance suppression is a universal property of admissible residue classes across all primorial scales and all constellation lengths $k \geq 2$.

\section{Extension to Prime Quads and Sextuplets: The Modulo-60 and Modulo-210 Lattices}

The framework generalizes to higher-order constellations. For prime quadruplets (e.g., $(p, p+2, p+6, p+8)$) and the associated sextuplet pattern, the modular constraints tighten further.

\subsection{Prime Quadruplets and Modulo-210}

A prime quadruplet $(p, p+2, p+6, p+8)$ requires $p$ to satisfy simultaneous congruences modulo 2, 3, and 5. The anchor $p$ must satisfy:
\begin{itemize}
  \item $p$ odd (mod 2: 1 choice),
  \item $p \equiv 2 \pmod 3$ (mod 3: 1 choice, since $p \equiv 1$ would give $p+2 \equiv 0 \pmod 3$ and $p+8 \equiv 0 \pmod 3$, hitting forbidden residues),
  \item $p \equiv 1 \pmod 5$ (mod 5: 1 choice, since $p+2, p+6, p+8$ cover 3 of the remaining 4 residues).
\end{itemize}
Thus at $P_3 = 30$, there is exactly 1 quadruplet-admissible residue class. The count grows as:
\begin{equation}
    A_{n,4} = \prod_{7 \le p \le p_n} (p-4),
\end{equation}
where the per-prime factor is $(p-4)$ because 4 of the $p$ residue classes are forbidden by the four positions of the quadruplet.

\begin{table}[h]
\centering
\caption{Admissible class counts for $k=4$ constellations}
\label{tab:quad_counts}
\begin{tabular}{cccc}
\toprule
$n$ & $P_n$ & Quad classes $A_{n,4}$ & $A_{n,4}/A_{n,3}$ \\
\midrule
3 & 30 & 1 & 0.50 \\
4 & 210 & 3 & 0.375 \\
5 & 2{,}310 & 21 & 0.328 \\
6 & 30{,}030 & 189 & 0.295 \\
\bottomrule
\end{tabular}
\end{table}

The ratio $A_{n,4}/A_{n,3}$ decreases as $(p-4)/(p-3) < 1$ for each prime $p \geq 7$, demonstrating that the combinatorial thinning accelerates with constellation order.

\section{Discussion: The Geometry of Variance Suppression}

In Paper 2, we observed that twin primes exhibit sub-Poisson variance across admissible classes. In this paper, we observe that prime triplets exhibit qualitatively the same effect. The continuous Hardy--Littlewood framework predicts asymptotic Poisson variance because it models prime events as independent. However, at any finite primorial scale, the discrete information geometry reveals a structured departure: cross-class counts have variance strictly below the Poisson limit.

This departure is strongest at small $n$, where the small number of admissible classes constrains the distribution of a fixed total count across bins---a familiar finite-population effect. As the primorial scale grows and the number of admissible classes explodes (from 4 at $n=3$ to 17{,}920 at $n=7$ for triplets), the structural boundary penalties relax systematically, and the CV/CV$_{\rm Poisson}$ ratio converges toward $1.0$. The discrete lattice boundary effects smoothly dissolve into continuous Poisson randomness at scale---a unifying picture across both twin primes ($k=2$) and triplets ($k=3$).

\section{Conclusion}

We have formalized the discrete information geometry for prime triplets, establishing the exact Fisher information on the modulo-30 lattice. By extending the empirical equidistribution census to $k=3$ constellations, we have demonstrated that hyper-equidistribution and variance suppression are universal features of prime constellations bounded by primorial lattices. 

This completes the foundational post-v4.3 RAPF trilogy. Paper 1 established the formal necessity of discrete lattice geometries and derived the exact Fisher information metric on the modulo-6 lattice. Paper 2 empirically verified uniform occupancy and sub-Poisson variance for twin primes across six primorial scales into $10^{14}$. This work extends the architecture to higher-order constellations, demonstrating that the sub-Poisson boundary effect observed for twins generalizes to prime triplets on the modulo-30 lattice, and that the suppression systematically relaxes toward Poisson as $n \to \infty$. Together, the three papers establish that prime $k$-tuples distribute with structured anti-clumping on primorial lattices at finite scales, with boundary effects that smoothly dissolve into the continuous Hardy--Littlewood picture at large scales.

\end{document}
